\documentclass[a4paper,11pt, tikz]{article}
\usepackage[T1]{fontenc}
\usepackage[utf8]{inputenc}
\usepackage{lmodern}
\usepackage{xcolor}
\usepackage{textcomp}
%\usepackage{newspaper}
%\date{\today}
%\currentvolume{1430}
%\currentissue{6}
%\usepackage{times}
\usepackage{graphicx}
\usepackage{multicol}
\usepackage{etoolbox}
\usepackage{picinpar}
%uasage of picinpar:
%\begin{window}[1,l,\includegraphics{},caption]xxxxx\end{window}
\usepackage{times}
\usepackage[T1]{fontenc}
\usepackage[utf8]{inputenc}
%\usepackage[nohead, nomarginpar, margin=1in, foot=.25in]{geometry}
\usepackage{lmodern}
%\usepackage{fontspec}
\usepackage{setspace}
%\usepackage{fancyhdr}
\usepackage{everyshi}
\usepackage{pdfpages}
\usepackage{indentfirst}
\usepackage[hyphens]{url}
\usepackage{hyperref}
\usepackage{pdfpages}
\usepackage{pgfpages}
\usepackage{amssymb}
\usepackage{amsmath}
\usepackage[notextcomp]{stix}
 \usepackage{paracol}
 \usepackage{stmaryrd}
 \usepackage{xurl}
 \usepackage{wasysym}
\usepackage{tikz}
 \usetikzlibrary{cd}
\usepackage[left=2cm, right=2cm, top=2cm, bottom=2cm]{geometry}
%\pagestyle{fancy}
%\lhead{Curriculum vitae, Anamitro Biswas}
%\chead{}
%\rhead{}
%\lfoot{}
%\cfoot{\thepage}
%\rfoot{}
%\makeatletter
%\renewcommand{\@chapapp}{Document}
%\makeatother
%\renewcommand\refname{\enhead\color{red}{References:}}
%\renewcommand{\headrulewidth}{0.4pt}
%\renewcommand{\footrulewidth}{0.4pt}
%\newfontfamily\bn{TiroBangla-Regular.ttf}[Script=Bengali]
%\newfontfamily\enhead{RubikDoodleShadow-Regular.ttf}[Script=Latin]
%\setmainfont{cmunrm.ttf}[Script=Latin]
%\newfontfamily\bn{TiroBangla-Regular.ttf}[Script=Bengali]
%\newfontfamily\latinfont{OfenbacherSchwabCAT.ttf}[Script=Latin]

%ALGEBRAIC TOPOLOGY:
\newcommand{\rel}{\text{ rel}}

%BRACKETS:
\newcommand{\lb}{\left(}
\newcommand{\rb}{\right)}
\newcommand{\lc}{\left\lbrace}
\newcommand{\rc}{\right\rbrace}
\newcommand{\ls}{\left[}
\newcommand{\rs}{\right]}

%CATEGORY THEORY:
\newcommand{\Ab}{\mathcal{A}\mathfrak{b}}
\newcommand{\Top}{\mathcal{T}\mathfrak{o}\mathfrak{p}}

%DASH AND BAR:
\newcommand{\dash}{^\prime}
\newcommand{\ddash}{^{\prime\prime}}
\newcommand{\tdash}{^{\prime\prime\prime}}

%FIELDS:
\newcommand{\reals}{\mathbb{R}}
\newcommand{\bC}{\mathbb{C}}
\newcommand{\bG}{\mathbb{G}}
\newcommand{\bP}{\mathbb{P}}
\newcommand{\bE}{\mathbb{E}}
\newcommand{\bH}{\mathbb{H}}
\newcommand{\bF}{\mathbb{F}}
\newcommand{\bN}{\mathbb{N}}
\newcommand{\naturals}{\mathbb{N}}
\newcommand{\Z}{\mathbb{Z}}

%FUZZY TOPOLOGY:
\newcommand{\nphi}{\bigodot}

%GREEK UPPERCASE:
\newcommand{\Lm}{\Lambda}
\newcommand{\G}{\Gamma}
\newcommand{\D}{\Delta}
\newcommand{\Th}{\Theta}
\newcommand{\F}{\Phi}

%GREEK LOWERCASE:
\newcommand{\m}{\mu}
\newcommand{\lm}{\lambda}
\newcommand{\vs}{\varsigma}
\newcommand{\g}{\gamma}
\renewcommand{\t}{\tau}
\renewcommand{\a}{\alpha}
\renewcommand{\b}{\beta}
\newcommand{\e}{\varepsilon}
\newcommand{\dt}{\delta}
\newcommand{\f}{\varphi}
\newcommand{\z}{\zeta}
\newcommand{\s}{\sigma}
\newcommand{\p}{\pi}
\newcommand{\y}{\eta}
\renewcommand{\th}{\theta}
\newcommand{\dl}{\partial}

%GROUP THEORY:
\newcommand{\inv}{^{-1}}

%HOMOLOGICAL ALGEBRA:
\newcommand{\Hom}{\text{Hom}}
\newcommand{\Ext}{\text{Ext}}
\newcommand{\im}{\text{im~}}
%\newcommand{\lsft}{shift left=1.5ext}

%HOMTOPY GROUPS:
\newcommand{\htopo}{\mathcal{H}\mathfrak{o}\mathcal{T}\mathfrak{o}\mathfrak{p}_*}
\newcommand{\topo}{\mathcal{T}\mathfrak{o}\mathfrak{p}_*}
\newcommand{\htop}{\mathcal{H}\mathfrak{o}\mathcal{T}\mathfrak{o}\mathfrak{p}}
\newcommand{\cattop}{\mathcal{T}\mathfrak{o}\mathfrak{p}}

%LOGIC:
\renewcommand{\implies}{\Rightarrow}

%MATHCAL:
\newcommand{\A}{\mathcal{A}}
\newcommand{\B}{\mathcal{B}}
\newcommand{\C}{\mathcal{C}}
\newcommand{\cD}{\mathcal{D}}
\newcommand{\E}{\mathcal{E}}
\newcommand{\cF}{\mathcal{F}}
\newcommand{\cG}{\mathcal{G}}
\newcommand{\Hy}{\mathcal{H}}
\newcommand{\I}{\mathcal{I}}
\newcommand{\J}{\mathcal{J}}
\newcommand{\List}{\mathcal{L}}
\newcommand{\M}{\mathcal{M}}
\newcommand{\N}{\mathcal{N}}
\newcommand{\R}{\mathcal{R}}
\newcommand{\Pa}{\mathcal{P}}
\newcommand{\La}{\mathcal{L}}
\newcommand{\Ta}{\mathcal{T}}
\newcommand{\mcS}{\mathcal{S}}
\newcommand{\V}{\mathcal{V}}
\newcommand{\X}{\mathcal{X}}
\newcommand{\Y}{\mathcal{Y}}

%MATHFRAK:
\newcommand{\Mf}{\mathfrak{M}}
\newcommand{\Nf}{\mathfrak{N}}
\newcommand{\gf}{\mathfrak{g}}
\newcommand{\mfi}{\mathfrak{i}}
\newcommand{\mfj}{\mathfrak{j}}
\newcommand{\mfa}{\mathfrak{A}}

%MATH BOLD:
\newcommand{\bfa}{\mathbf{a}}
\newcommand{\bfb}{\mathbf{b}}
\newcommand{\bfg}{\mathbf{g}}
\newcommand{\bft}{\mathbf{t}}
\newcommand{\bfx}{\mathbf{x}}
\newcommand{\bfy}{\mathbf{y}}
\newcommand{\bfzero}{\mathbf{0}}


%MAPS:
\newcommand{\tends}{\rightarrow}
\newcommand{\id}{\text{id}}



%SET THEORY:
\newcommand{\cl}{\overline}
\newcommand{\intr}{^\circ}
\newcommand{\ir}{^\text{int}}
\newcommand{\comp}{^C}

%HEADINGS:
\newtheorem{thm}{Theorem}
\newtheorem{prop}[thm]{Proposition}
\newtheorem{cor}[thm]{Corollary}
\newtheorem{defn}[thm]{Definition}
\newtheorem{obs}[thm]{Observation}
%\newtheorem{prop}[thm]{Proposition}
\newtheorem{lem}[thm]{Lemma}
\newtheorem{conj}[thm]{Conjecture}
\newtheorem{claim}[thm]{Claim}
\newtheorem{note}[thm]{Note}
\newtheorem{notation}[thm]{Notation}
\newtheorem{rem}[thm]{Remark}
\newtheorem{prob}[thm]{Problem}
\newcommand{\Ob}{\text{Ob}}
\newcommand{\Eq}{\text{Eq}}

%\setmainfont{Cormorant Garamond SemiBold}
\urlstyle{same}
\usepackage{tikzpagenodes}
\usepackage{tikz}
\usepackage{eso-pic}
\begin{document}
\title{Ext Functor}
\author{Anamitro Biswas\\Email: anamitroappu@gmail.com}
\maketitle
\section{Projective Resolution}
\noindent{}Let $P$ be a module. For any surjective map $g:B\tends C\tends 0$. if $f:P\tends C$ always has a lifting $h:P\tends B$, we say that $P$ is \emph{projective}.
Let $A$ be an abelian group. Consider a projective resolution of $A$
$$P_\#: \dots\xrightarrow{\dl_{n+2}} P_{n+1}\xrightarrow{\dl_{n+1}} P_n\xrightarrow{\dl_{n}} P_{n-1}\xrightarrow{\dl_{n-1}}\dots\xrightarrow{\dl_{3}} P_2\xrightarrow{\dl_{2}} P_1\xrightarrow{\dl_{1}} P_0\xrightarrow{\dl_{0}} A\xrightarrow{\eta} 0,$$
where each $P_i$ is projective and the sequence is exact. Additionally, if each $P_i$ is free, we call $P_\#$ a \emph{free resolution}.
For a commutative ring with $1$, $R$, applying the $\Hom_R$ functor with an abelian group $G$, we get a cochain complex that might not be exact:
$$0\xrightarrow{d_{\y}}\Hom_R\lb A, G\rb\xrightarrow{d_{0}}\Hom_R\lb P_0, G\rb\xrightarrow{d_{1}}\Hom_R\lb P_1, G\rb\xrightarrow{d_{2}}\dots$$
The $n$th cohomology group of this cochain complex gives sort of a measurement of the non-exactness of this cochain complex at the $n$th position. We define that $n$th cohomology group as the $\Ext$ functor from $\Ab^2$ to $\Ab$,
$$\Ext_R^n\lb A, G\rb\coloneq \ker d_{n+1}/\im d_n.$$

It can be proved that this $\Ext$ functor does not depend on the particular projective resolution taken, and is just a function of $A$ and $G$.

If the cochain complex is exact at any point, which will be the case if $A$ is a free abelian group and $P$ a free resolution of $A$, then $\Ext_R^n\lb A, G\rb$ is trivial there.\\

\noindent{}\textit{If $A$ is a free abelian group.} Then the free abelian group generated by $A$ is $A$. Thus we have the chain $0\xrightarrow{} R=\ker\f\xrightarrow{} \cF(A)=A\xrightarrow{\f} A\xrightarrow{} 0$, where $\f$ ends up being the identity map, and hence $R=0$. So $\dl_0$ is just the identity map on $A$, and $\dl_1$ takes whole of $A$ to $0$; accordingly, $\Ext_R^0\lb A, G\rb\coloneq \ker d_{1}/\im d_0$ is trivial. It is also conversely true that a trivial $\Ext^0$ means that the concerned entry abelian group is free. If we consider the exact sequence with the canonical homomorphism $\cF(A)\tends A$, by assumption it splits which implies that $A$ is projective (because it's a direct summand of the free abelian group $\cF(A)$), hence free, abelian group.\\

\section{Examples}
\noindent{}\textbf{Example 1.} Let $A=\Z/m\Z$. Consider
\begin{equation*}
\begin{tikzcd}
0 \arrow[r, "\dl_2"] & \Z \arrow[r, "\dl_1", "m"'] & \Z \arrow[r, "\dl_0"] & {\Z/m\Z} \arrow[r] & 0
\end{tikzcd}
\end{equation*}
This gives us (taking $R=\Z$)
\begin{equation*}
\begin{tikzcd}
0 \arrow[r] & \Hom\lb \Z/m\Z, G\rb\arrow[r, "d_0"] & \Hom\lb \Z, G\rb \arrow[r, "d_1"] & \Hom\lb \Z, G\rb \arrow[r, "d_2"] & 0
\end{tikzcd}
\end{equation*}
We have the following commutative diagram:
\begin{equation*}
\begin{tikzcd}
&  &\\
\color{red}{\cl{m\b}\in}\color{black}{}\Z/m\Z \arrow[dd, "\f"] & \Z\color{red}{\ni m\b}\color{black}{} \arrow[l, "\p"'] \arrow [ddl, dashrightarrow, "\psi"] & \\
& & \Z\color{red}{\ni\b}\color{black}{} \arrow[ul, "m"] \arrow[dll, dashrightarrow, "m\psi"]\\
\color{red}{g\in}\color{black}{} G & &
\end{tikzcd}
\end{equation*}
Here, $d_0\lb \f\rb=\f\circ \pi$, hence $\im d_0=\lc \text{periodic functions with period } m\rc$ and $\ker d_1=\lc \psi: \Z\tends G~\mid~\psi\lb m\b\rb\right.$ $\left.=0~\forall~\b\in\Z\rc$. Therefore
$$\Ext^0\lb \Z/m\Z, G\rb=\dfrac{\ker d_1}{\im d_0}=\dfrac{\lc\psi:\Z\tends G~|~\psi\lb m\b\rb=m\psi\lb \b\rb=0~\forall~\b\in\Z\rc}{\lc \psi:\Z\tends G~|~\psi\lb \a+m\rb=\psi\lb \a\rb~\forall~\a\in\Z\rc}.$$
Now, for a $[\psi]\in\Ext^0\lb \Z/m\Z, G\rb$, $m\psi\equiv 0\implies m\psi\lb \b\rb=0~\forall~\b\in\Z\implies m\psi\lb \b+m\rb=0\implies m\psi(\b)+m\psi(m)=0\text{ since $\psi$ is a homomorphism }\implies m\psi(m)=0$.

Define $\z:\Hom\lb \Z, G\rb\tends G;~\psi\mapsto \psi(m)=g\in G$. This is well-defined. Therefore I can identify classes of $\psi$ with corresponding elements of $g$ such that $mg=0_G$. So, $\Ext^0\lb \Z/m\Z, G\rb\approx {}_mG=\lc g\in G~|~mg=0\rc$.

Now, let us calculate $\Ext^1\lb \Z/m\Z, G\rb=\dfrac{\ker d_2}{\im d_1}=\dfrac{\Hom\lb \Z, G\rb}{m\Hom\lb \Z, G\rb}$.\\

\begin{equation*}
\begin{tikzcd}
\Z \arrow[dr, dashrightarrow, "d_1\f"] &&  \Z \arrow[ll, "m"']\arrow[dl, "\f"]\\
&G&
\end{tikzcd}
\end{equation*}
If $\f\in\Hom\lb\Z, G\rb$, as soon as we define $\f(1)$, $\f$ is fully defined. So, one $\f$ is equivalent to one element $g\in G$ where $\f(1)$ is assigned. So, $\Hom\lb\Z, G\rb\approx G\implies Ext^1\lb \Z/m\Z, G\rb\approx G/m G$.

Note here that if $|G|\not\mid m$, then $\Ext^0\lb \Z/m\Z, G\rb$ is non-trivial, providing us with an example of a non-trivial $\Hom$-cochain and $\Ext$, showing that the definition is vacuously true. Furthermore, $\Hom\lb\Z/m\Z, G\rb\tends \Hom\lb\Z, G\rb$ is non-trivial only if $G$ has a torsion.
%\textbf{Example 2.}

\section{Independence of $\Ext$ on resolution}

\begin{thm}
Let $A, A\dash\in\Ab$. If a homomorphism $f:A\tends A\dash$ exists, it would induce a chain map $f_\#$ in the following scenario:

\begin{equation*}
\begin{tikzcd}
\dots\ar[r, "\dl_3"] &P_2\arrow[r, "\dl_2"] \ar[d, dashrightarrow, "f_2"] &P_1\arrow[r, "\dl_1"] \ar[d, dashrightarrow, "f_1"]&P_0\arrow[r, "\dl_0"] \ar[d, dashrightarrow, "f_0"]& A\ar[r, "\y"] \ar[d, "f"]& 0\\
\dots\ar[r, "\dl_3"] &P_2\dash\arrow[r, "\dl_2\dash"]  &P_1\dash\arrow[r, "\dl_1\dash"]&P_0\dash\arrow[r, "\dl_0\dash"] & A\dash\ar[r, "\y\dash"]& 0
\end{tikzcd}
\end{equation*}
if each $P_i$ in the above row is projective and the bottom row is exact. Moreover, any two such chain maps will be homotopic.
\end{thm}
We prove by induction on $n\geq 0$. Since $P_0$ is projective, the following commutative diagram proves the existence of $f_0$
\begin{equation*}
\begin{tikzcd}
& P_0\ar[dl, dashrightarrow, "f_0"'] \ar[d, "f\dl_0"]\\
P_0\dash\ar[r, "\dl_0\dash"']& A\dash
\end{tikzcd}
\end{equation*}
For the induction step, consider
\begin{equation*}
\begin{tikzcd}
P_{n+1}\ar[r, "\dl_{n+1}"]\ar[d, dashrightarrow] & P_n\ar[r, "\dl_n"]\ar[d, "f_n"] & P_{n-1}\ar[r]\ar[d, "f_{n-1}"] & \dots\\
P_{n+1}\dash\ar[r, "\dl_{n+1}\dash"] & P_n\dash\ar[r, "\dl_n\dash"] & P_{n-1}\dash\ar[r] & \dots
\end{tikzcd}
\end{equation*}
We can have this implication due to projectivity of $P_{n+1}$
\begin{equation*}
\begin{tikzcd}
& P_{n+1}\ar[dl, dashrightarrow, "f_{n+1}"']\ar[d, "f_n\dl_{n+1}"]&\\
P_{n+1}\dash\ar[r, "\dl_{n+1}\dash"'] & \im\dl_{n+1}\dash\ar [r] & 0
\end{tikzcd}
\end{equation*}
only if $\im\lb f_n\dl_{n+1}\rb\subset \im \dl_{n-1}\dash=\ker\dl_n\dash$, which is true since $\dl_n\dash f_n\dl_{n+1}=f_{n-1}\dl_n\dl_{n+1}=f_{n-1}0=0$. Thus by induction we have proved the existence of the chain map $f_\#$. On to its uniqueness up to homotopy equivalence.

Consider the chain complexes where we have to construct the homotopy map $s_\#$,
\begin{equation*}
\begin{tikzcd}
{} &P_{n+1}\ar[r, "\dl_{n+1}"]\ar[d, "h_{n+1}"', shift right=3ex, blue] \ar[d, "f_{n+1}"', shift left=1ex, black] \ar[dl, "s_{n+1}"', red, shift right =2.5ex, dashrightarrow] & P_n\ar[r, "\dl_n"]\ar[d, "h_{n}"', shift right=1.5ex, blue] \ar[d, "f_{n}"', shift left=1.5ex, black] \ar[dl, "s_{n}"', red]  & P_{n-1}\ar[r]\ar[d, "h_{n-1}"', shift right=2ex, blue] \ar[dl, "s_{n-1}"', red] \ar[d, "f_{n-1}"', shift left=2ex, black]  & \dots \ar[r, "\dl_3"] &P_2\arrow[r, "\dl_2"] \ar[d, "h_{2}"', shift right=1.5ex, blue] \ar[d, "f_{2}"', shift left=1.5ex, black] \ar[dl, "s_{2}"', red] &P_1\arrow[r, "\dl_1"] \ar[d, "h_{1}"', shift right=1.5ex, blue] \ar[dl, "s_{1}"', red] \ar[d, "f_{1}"', shift left=1.5ex, black] &P_0\arrow[r, "\dl_0"] \ar[d, "h_{0}"', shift right=1.5ex, blue] \ar[dl, "s_{0}"', red]\ar[d, "f_{0}"', shift left=1.5ex, black] & A\ar[r, "\y"] \ar[dl, "s_{-1}"', red] \ar[d, "f"]& 0 \ar[dl, "s_{-2}"', red]\\
{} &P_{n+1}\dash\ar[r, "\dl_{n+1}\dash"] & P_n\dash\ar[r, "\dl_n\dash"] & P_{n-1}\dash\ar[r] & \dots \ar[r, "\dl_3"] &P_2\dash\arrow[r, "\dl_2\dash"]  &P_1\dash\arrow[r, "\dl_1\dash"]&P_0\arrow[r, "\dl_0\dash"] & A\dash\ar[r, "\y\dash"]& 0
\end{tikzcd}
\end{equation*}
such that $s_{n-1}\dl_n+\dl_{n+1}s_n=h_n-f_n~\forall~n\in\naturals\cup\lc -2, -1, 0\rc$. Putting $s_{-1}=0$ gives this for $n=0$. As for the inductive step, we have to show that the $s_n$'s so constructed satisfies $h_{n+1}-f_{n+1}-s_n\dl_{n+1}=\dl_{n+2}\dash s_{n+1}$ for the $(n+1)$th level. We have that $\lb h_{n+1}-f_{n+1}-s_n\dl_{n+1}\rb\lb P_{n+1}\rb\subset P_{n+1}\dash$. Only if we can show that $\lb h_{n+1}-f_{n+1}-s_n\dl_{n+1}\rb\lb P_{n+1}\rb\subset \im\dl_{n+2}\dash=\ker\dl_{n+1}\dash$, then by surjectivity of $\dl_{n+2}\dash$ on its codomain, we have the following commutative diagram from where the projectivity of $P_{n+1}$ implies the existence of $s_{n+1}$ with the desired property.
\begin{equation*}
\begin{tikzcd}
& P_{n+1}\ar[dl, dashrightarrow, "s_{n+1}"'] \ar[d, "h_{n+1}-f_{n+1}-s_n\dl_{n+1}"] &\\
P_{n+2}\dash\ar[r, "\dl_{n+2}"] & P_{n+1}\dash\ar[r]& 0
\end{tikzcd}
\end{equation*}
\begin{eqnarray*}
 \text{But } \dl_{n+1}\dash\lb h_{n+1}-f_{n+1}-s_n\dl_{n+1}\rb &=& \dl_{n+1}\dash\lb h_{n+1}-f_{n+1}\rb-\lb \dl_{n+1}\dash s_n\rb\dl_{n+1}\\
  &=& \dl_{n+1}\dash\lb h_{n+1}-f_{n+1}\rb-\lb h_n-f_n-s_{n-1}\dl_n\rb\dl_{n+1}~\ls\text{by induction hypothesis}\rs\\
  &=& 0,~\text{since $h$ and $f$ are chain maps}.\square 
\end{eqnarray*}
This being established, we try to see where this homotopy maps through the $\Hom$ functor. Of course, for $\f:P_n\dash\tends G$, we get a map $\f\circ h_n:P_n\tends G$, so $h_n$ induces $\th: \Hom\lb P_n\dash, G\rb\tends \Hom\lb P_n, G\rb; \f\mapsto \f\circ h_n$. In a similar manner all the chain and homotopy maps are transferred to our commutative diagram through $\Hom$
\begin{equation*}
\begin{tikzcd}
{\Hom\lb P_{n-1}, G\rb}\ar[r, "\D_n"] &{ \Hom\lb P_n, G\rb }& \\
& {\Hom\lb P_n\dash, G\rb }\ar[ul, "\s_{n-1}"] \ar[u, "\th_n", shift left=1.5ex] \ar[u, "\f_n", shift right=1.5ex] \ar[r, "\D_{n+1}\dash"] &{ \Hom\lb P_{n+1}\dash, G\rb} \ar[ul, "\s_n"]
\end{tikzcd}
\end{equation*}
In fact, a cochain complex is induced. As we know, $\Hom$ is a contravariant functor, and as shown below, the chain homotopy condition obtained previously gives $\th_n-\f_n=\s_n\D_{n+1}\dash+\D_n\s_{n-1}$, which gives the cochain homotopy relation between the two cochains obtained through $\Hom$ functor.

Now, if $A=A\dash$, $f$ is instead an isomorphism, and the bottom row is projective as well, i.e., we are dealing with two projective resolutions of $A$. An ismorphism $f$ induces isomorphisms $P_n\tends P_n\dash$ by inducing homomorphisms both ways. All induced chain maps have been seen to be homtopic to the isomorphism, i.e., resolutions $P_\#$ and $P_\#\dash$ are homotopically equivalent. We know that cohomology groups of homotopic cochain maps are isomorphic. So the $\Ext$ functors, which are cohomology groups of the $\Hom$ cochains of the resolutions, are the same. That is, they are independent of the particular resolution chosen over $A$, and instead just depend on $A$ and $G$-- but also on the ring $R$ over which the homomorphisms are taken.
\end{document}
